\documentclass[10pt,a4paper]{amsart}
\usepackage{amsmath,amssymb,amsthm,graphicx}
\usepackage{verbatim}
\newenvironment{reference}{\comment}{\endcomment}
%\newenvironment{reference}{}{}
\newenvironment{slogan}{\comment}{\endcomment}
\newenvironment{history}{\comment}{\endcomment}

% For commutative diagrams we use Xy-pic
\usepackage[all]{xy}

% We use 2cell for 2-commutative diagrams.
\xyoption{2cell}
% Theorem environments.
%
\theoremstyle{plain}
\newtheorem{theorem}[subsection]{Theorem}
\newtheorem{proposition}[subsection]{Proposition}
\newtheorem{lemma}[subsection]{Lemma}

\theoremstyle{definition}
\newtheorem{definition}[subsection]{Definition}
\newtheorem{example}[subsection]{Example}
\newtheorem{exercise}[subsection]{Exercise}
\newtheorem{situation}[subsection]{Situation}

\theoremstyle{remark}
\newtheorem{remark}[subsection]{Remark}
\newtheorem{remarks}[subsection]{Remarks}

\numberwithin{equation}{subsection}

% Macros
%
\def\lim{\mathop{\mathrm{lim}}\nolimits}
\def\colim{\mathop{\mathrm{colim}}\nolimits}
\def\Spec{\mathop{\mathrm{Spec}}}
\def\Hom{\mathop{\mathrm{Hom}}\nolimits}
\def\Ext{\mathop{\mathrm{Ext}}\nolimits}
\def\SheafHom{\mathop{\mathcal{H}\!\mathit{om}}\nolimits}
\def\SheafExt{\mathop{\mathcal{E}\!\mathit{xt}}\nolimits}
\def\Sch{\mathit{Sch}}
\def\Mor{\mathop{\mathrm{Mor}}\nolimits}
\def\Ob{\mathop{\mathrm{Ob}}\nolimits}
\def\Sh{\mathop{\mathit{Sh}}\nolimits}
\def\NL{\mathop{N\!L}\nolimits}
\def\CH{\mathop{\mathrm{CH}}\nolimits}
\def\proetale{{pro\text{-}\acute{e}tale}}
\def\etale{{\acute{e}tale}}
\def\QCoh{\mathit{QCoh}}
\def\Ker{\mathop{\mathrm{Ker}}}
\def\Im{\mathop{\mathrm{Im}}}
\def\Coker{\mathop{\mathrm{Coker}}}
\def\Coim{\mathop{\mathrm{Coim}}}

% Boxtimes
%
\DeclareMathSymbol{\boxtimes}{\mathbin}{AMSa}{"02}

%
% Macros for moduli stacks/spaces
%
\def\QCohstack{\mathcal{QC}\!\mathit{oh}}
\def\Cohstack{\mathcal{C}\!\mathit{oh}}
\def\Spacesstack{\mathcal{S}\!\mathit{paces}}
\def\Quotfunctor{\mathrm{Quot}}
\def\Hilbfunctor{\mathrm{Hilb}}
\def\Curvesstack{\mathcal{C}\!\mathit{urves}}
\def\Polarizedstack{\mathcal{P}\!\mathit{olarized}}
\def\Complexesstack{\mathcal{C}\!\mathit{omplexes}}
% \Pic is the operator that assigns to X its picard group, usage \Pic(X)
% \Picardstack_{X/B} denotes the Picard stack of X over B
% \Picardfunctor_{X/B} denotes the Picard functor of X over B
\def\Pic{\mathop{\mathrm{Pic}}\nolimits}
\def\Picardstack{\mathcal{P}\!\mathit{ic}}
\def\Picardfunctor{\mathrm{Pic}}
\def\Deformationcategory{\mathcal{D}\!\mathit{ef}}
\usepackage{cleveref}
\crefname{originallemma}{original lemma}{original lemmas}
\Crefname{originallemma}{Original lemma}{Original lemmas}
\makeatletter
\newlabel{chow-lemma-cycle-push-sheaf}{{02R6}{source}{}{sourcepointer-1}{}}
\newlabel{chow-lemma-cycle-push-sheaf@cref}{{[originallemma][0][]02R6}{[1][0][]source}{}{}{}}
\makeatother
\crefname{originalreference}{original reference}{original references}
\Crefname{originalreference}{Original reference}{Original references}
\makeatletter
\newlabel{lemma-cycle-closed-coherent}{{02QY}{source}{}{sourcepointer-2}{}}
\newlabel{lemma-cycle-closed-coherent@cref}{{[originalreference][0][]02QY}{[1][0][]source}{}{}{}}
\makeatother
\crefname{originalreference}{original reference}{original references}
\Crefname{originalreference}{Original reference}{Original references}
\makeatletter
\newlabel{lemma-length-closed-immersion}{{0EE3}{source}{}{sourcepointer-3}{}}
\newlabel{lemma-length-closed-immersion@cref}{{[originalreference][0][]0EE3}{[1][0][]source}{}{}{}}
\makeatother
\crefname{originalreference}{original reference}{original references}
\Crefname{originalreference}{Original reference}{Original references}
\makeatletter
\newlabel{lemma-point-of-max-dimension}{{0EE6}{source}{}{sourcepointer-4}{}}
\newlabel{lemma-point-of-max-dimension@cref}{{[originalreference][0][]0EE6}{[1][0][]source}{}{}{}}
\makeatother
\crefname{originalreference}{original reference}{original references}
\Crefname{originalreference}{Original reference}{Original references}
\makeatletter
\newlabel{situation-setup}{{02QL}{source}{}{sourcepointer-5}{}}
\newlabel{situation-setup@cref}{{[originalreference][0][]02QL}{[1][0][]source}{}{}{}}
\makeatother
\crefname{originallemma}{original lemma}{original lemmas}
\Crefname{originallemma}{Original lemma}{Original lemmas}
\makeatletter
\newlabel{spaces-cohomology-lemma-coherent-support-closed}{{07UG}{source}{}{sourcepointer-6}{}}
\newlabel{spaces-cohomology-lemma-coherent-support-closed@cref}{{[originallemma][0][]07UG}{[1][0][]source}{}{}{}}
\makeatother
\crefname{originallemma}{original lemma}{original lemmas}
\Crefname{originallemma}{Original lemma}{Original lemmas}
\makeatletter
\newlabel{spaces-cohomology-lemma-proper-pushforward-coherent}{{08AR}{source}{}{sourcepointer-7}{}}
\newlabel{spaces-cohomology-lemma-proper-pushforward-coherent@cref}{{[originallemma][0][]08AR}{[1][0][]source}{}{}{}}
\makeatother
\crefname{originaldefinition}{original definition}{original definitions}
\Crefname{originaldefinition}{Original definition}{Original definitions}
\makeatletter
\newlabel{spaces-morphisms-definition-scheme-theoretic-image}{{082Y}{source}{}{sourcepointer-8}{}}
\newlabel{spaces-morphisms-definition-scheme-theoretic-image@cref}{{[originaldefinition][0][]082Y}{[1][0][]source}{}{}{}}
\makeatother
\crefname{originallemma}{original lemma}{original lemmas}
\Crefname{originallemma}{Original lemma}{Original lemmas}
\makeatletter
\newlabel{spaces-morphisms-lemma-base-change-proper}{{04WP}{source}{}{sourcepointer-9}{}}
\newlabel{spaces-morphisms-lemma-base-change-proper@cref}{{[originallemma][0][]04WP}{[1][0][]source}{}{}{}}
\makeatother
\crefname{originallemma}{original lemma}{original lemmas}
\Crefname{originallemma}{Original lemma}{Original lemmas}
\makeatletter
\newlabel{spaces-morphisms-lemma-closed-immersion-proper}{{04XZ}{source}{}{sourcepointer-10}{}}
\newlabel{spaces-morphisms-lemma-closed-immersion-proper@cref}{{[originallemma][0][]04XZ}{[1][0][]source}{}{}{}}
\makeatother
\crefname{originallemma}{original lemma}{original lemmas}
\Crefname{originallemma}{Original lemma}{Original lemmas}
\makeatletter
\newlabel{spaces-morphisms-lemma-quasi-compact-scheme-theoretic-image}{{082Z}{source}{}{sourcepointer-11}{}}
\newlabel{spaces-morphisms-lemma-quasi-compact-scheme-theoretic-image@cref}{{[originallemma][0][]082Z}{[1][0][]source}{}{}{}}
\makeatother
\crefname{originallemma}{original lemma}{original lemmas}
\Crefname{originallemma}{Original lemma}{Original lemmas}
\makeatletter
\newlabel{spaces-morphisms-lemma-universally-closed-permanence}{{04NX}{source}{}{sourcepointer-12}{}}
\newlabel{spaces-morphisms-lemma-universally-closed-permanence@cref}{{[originallemma][0][]04NX}{[1][0][]source}{}{}{}}
\makeatother
\crefname{originallemma}{original lemma}{original lemmas}
\Crefname{originallemma}{Original lemma}{Original lemmas}
\makeatletter
\newlabel{spaces-over-fields-lemma-finite-in-codim-1}{{0AD2}{source}{}{sourcepointer-13}{}}
\newlabel{spaces-over-fields-lemma-finite-in-codim-1@cref}{{[originallemma][0][]0AD2}{[1][0][]source}{}{}{}}
\makeatother
\crefname{originalproposition}{original proposition}{original propositions}
\Crefname{originalproposition}{Original proposition}{Original propositions}
\makeatletter
\newlabel{spaces-properties-proposition-locally-quasi-separated-open-dense-scheme}{{06NH}{source}{}{sourcepointer-14}{}}
\newlabel{spaces-properties-proposition-locally-quasi-separated-open-dense-scheme@cref}{{[originalproposition][0][]06NH}{[1][0][]source}{}{}{}}
\makeatother
\title{Proper pushforward of a coherent sheaf of delta-dimension at most k commutes with the k-dimensional cycle class}
\author{The Stacks Project Authors}
\date{stacks-project:spaces-chow.tex:local-snapshot-20261012; GNU FDL 1.2 or later}
\begin{document}
\maketitle
\noindent\textit{Editorial note.} All delivered mathematical text is the author's own, copied from the pinned Stacks Project chapter file spaces-chow.tex; no statement or proof step was rewritten, reordered or removed.\par\medskip
\noindent\textit{Editorial note.} The counted proof is the author's complete argument: a support-and-scheme-theoretic-image reduction to closed immersions followed by a component comparison and a shrink to the finite affine case\par\medskip
\noindent\textit{Editorial note.} The label lemma-cycle-push-sheaf names the counted statement, and the proof environment immediately following it in the source is delivered verbatim as its proof.\par\medskip
\noindent\textit{Editorial note.} The counted statement is carried by the author's own lemma environment in the chapter source.\par\medskip
\noindent\textit{Allowed standard prerequisites (editorial).}\begin{itemize}
\item Cohomology of Spaces, proper pushforward of a coherent module is coherent, and every coherent sheaf on a good algebraic space is the pushforward of a coherent module on a closed subspace that is its support
\item The closed immersion case of the cycle class computation for a pushforward, and the identification of the k-dimensional part of a cycle class with a sum over the irreducible components of delta-dimension k
\item Spaces over Fields, finiteness in codimension 1 near a generic point of a component, and Chow Homology's cycle pushforward lemma for schemes, used as the final reduction step
\end{itemize}\medskip
\ifdefined\hypertarget\hypertarget{sourcepointer-1}{}\fi
\noindent\textit{Source reference.} Original lemma 02R6: sources/stacks-chow/tags.txt, source lines 3613--3613. Prerequisite pointer invoked at proof line(s) 1117; it names a standard fact of a sibling chapter that the counted argument uses, and every step of the counted proof built on it is retained verbatim in the delivered body.\par\medskip
\ifdefined\hypertarget\hypertarget{sourcepointer-2}{}\fi
\noindent\textit{Source reference.} Original reference 02QY: sources/stacks-spaces-chow/pinned-original-source.tex, source lines 762--762. Prerequisite pointer invoked at proof line(s) 1043; it names a standard fact of the same chapter that the counted argument uses, and every step of the counted proof built on it is retained verbatim in the delivered body.\par\medskip
\ifdefined\hypertarget\hypertarget{sourcepointer-3}{}\fi
\noindent\textit{Source reference.} Original reference 0EE3: sources/stacks-spaces-chow/pinned-original-source.tex, source lines 521--521. Prerequisite pointer invoked at proof line(s) 1080; it names a standard fact of the same chapter that the counted argument uses, and every step of the counted proof built on it is retained verbatim in the delivered body.\par\medskip
\ifdefined\hypertarget\hypertarget{sourcepointer-4}{}\fi
\noindent\textit{Source reference.} Original reference 0EE6: sources/stacks-spaces-chow/pinned-original-source.tex, source lines 626--626. Prerequisite pointer invoked at proof line(s) 1099; it names a standard fact of the same chapter that the counted argument uses, and every step of the counted proof built on it is retained verbatim in the delivered body.\par\medskip
\ifdefined\hypertarget\hypertarget{sourcepointer-5}{}\fi
\noindent\textit{Source reference.} Original reference 02QL: sources/stacks-spaces-chow/pinned-original-source.tex, source lines 264--264. Prerequisite pointer invoked at proof line(s) 1022; it names a standard fact of the same chapter that the counted argument uses, and every step of the counted proof built on it is retained verbatim in the delivered body.\par\medskip
\ifdefined\hypertarget\hypertarget{sourcepointer-6}{}\fi
\noindent\textit{Source reference.} Original lemma 07UG: sources/stacks-chow/tags.txt, source lines 9971--9971. Prerequisite pointer invoked at proof line(s) 1047; it names a standard fact of a sibling chapter that the counted argument uses, and every step of the counted proof built on it is retained verbatim in the delivered body.\par\medskip
\ifdefined\hypertarget\hypertarget{sourcepointer-7}{}\fi
\noindent\textit{Source reference.} Original lemma 08AR: sources/stacks-chow/tags.txt, source lines 10541--10541. Prerequisite pointer invoked at proof line(s) 1039; it names a standard fact of a sibling chapter that the counted argument uses, and every step of the counted proof built on it is retained verbatim in the delivered body.\par\medskip
\ifdefined\hypertarget\hypertarget{sourcepointer-8}{}\fi
\noindent\textit{Source reference.} Original definition 082Y: sources/stacks-chow/tags.txt, source lines 10268--10268. Prerequisite pointer invoked at proof line(s) 1053; it names a standard fact of a sibling chapter that the counted argument uses, and every step of the counted proof built on it is retained verbatim in the delivered body.\par\medskip
\ifdefined\hypertarget\hypertarget{sourcepointer-9}{}\fi
\noindent\textit{Source reference.} Original lemma 04WP: sources/stacks-chow/tags.txt, source lines 6367--6367. Prerequisite pointer invoked at proof line(s) 1067; it names a standard fact of a sibling chapter that the counted argument uses, and every step of the counted proof built on it is retained verbatim in the delivered body.\par\medskip
\ifdefined\hypertarget\hypertarget{sourcepointer-10}{}\fi
\noindent\textit{Source reference.} Original lemma 04XZ: sources/stacks-chow/tags.txt, source lines 6412--6412. Prerequisite pointer invoked at proof line(s) 1068; it names a standard fact of a sibling chapter that the counted argument uses, and every step of the counted proof built on it is retained verbatim in the delivered body.\par\medskip
\ifdefined\hypertarget\hypertarget{sourcepointer-11}{}\fi
\noindent\textit{Source reference.} Original lemma 082Z: sources/stacks-chow/tags.txt, source lines 10269--10269. Prerequisite pointer invoked at proof line(s) 1064; it names a standard fact of a sibling chapter that the counted argument uses, and every step of the counted proof built on it is retained verbatim in the delivered body.\par\medskip
\ifdefined\hypertarget\hypertarget{sourcepointer-12}{}\fi
\noindent\textit{Source reference.} Original lemma 04NX: sources/stacks-chow/tags.txt, source lines 6095--6095. Prerequisite pointer invoked at proof line(s) 1069; it names a standard fact of a sibling chapter that the counted argument uses, and every step of the counted proof built on it is retained verbatim in the delivered body.\par\medskip
\ifdefined\hypertarget\hypertarget{sourcepointer-13}{}\fi
\noindent\textit{Source reference.} Original lemma 0AD2: sources/stacks-chow/tags.txt, source lines 13075--13075. Prerequisite pointer invoked at proof line(s) 1106; it names a standard fact of a sibling chapter that the counted argument uses, and every step of the counted proof built on it is retained verbatim in the delivered body.\par\medskip
\ifdefined\hypertarget\hypertarget{sourcepointer-14}{}\fi
\noindent\textit{Source reference.} Original proposition 06NH: sources/stacks-chow/tags.txt, source lines 8536--8536. Prerequisite pointer invoked at proof line(s) 1112; it names a standard fact of a sibling chapter that the counted argument uses, and every step of the counted proof built on it is retained verbatim in the delivered body.\par\medskip
\begin{lemma}
\label{lemma-cycle-push-sheaf}
In Situation \ref{situation-setup} let $f : X \to Y$ be a proper morphism
of good algebraic spaces over $B$.
\begin{enumerate}
\item Let $Z \subset X$ be a closed subspace with $\dim_\delta(Z) \leq k$.
Then
$$
f_*[Z]_k = [f_*{\mathcal O}_Z]_k.
$$
\item Let $\mathcal{F}$ be a coherent sheaf on $X$ such that
$\dim_\delta(\text{Supp}(\mathcal{F})) \leq k$. Then
$$
f_*[\mathcal{F}]_k = [f_*{\mathcal F}]_k.
$$
\end{enumerate}
Note that the statement makes sense since $f_*\mathcal{F}$ and
$f_*\mathcal{O}_Z$ are coherent $\mathcal{O}_Y$-modules by
Cohomology of Spaces, Lemma
\ref{spaces-cohomology-lemma-proper-pushforward-coherent}.
\end{lemma}

\begin{proof}
Part (1) follows from (2) and Lemma \ref{lemma-cycle-closed-coherent}.
Let $\mathcal{F}$ be a coherent sheaf on $X$.
Assume that $\dim_\delta(\text{Supp}(\mathcal{F})) \leq k$.
By Cohomology of Spaces, Lemma
\ref{spaces-cohomology-lemma-coherent-support-closed}
there exists a closed immersion $i : Z \to X$ and a coherent
$\mathcal{O}_Z$-module $\mathcal{G}$ such that
$i_*\mathcal{G} \cong \mathcal{F}$ and such that the support
of $\mathcal{F}$ is $Z$. Let $Z' \subset Y$ be the scheme theoretic image
of $f|_Z : Z \to Y$, see Morphisms of Spaces, Definition
\ref{spaces-morphisms-definition-scheme-theoretic-image}.
Consider the commutative diagram
$$
\xymatrix{
Z \ar[r]_i \ar[d]_{f|_Z} &
X \ar[d]^f \\
Z' \ar[r]^{i'} & Y
}
$$
of algebraic spaces over $B$. Observe that $f|_Z$ is surjective
(follows from Morphisms of Spaces, Lemma
\ref{spaces-morphisms-lemma-quasi-compact-scheme-theoretic-image}
and the fact that $|f|$ is closed) and proper
(follows from Morphisms of Spaces, Lemmas
\ref{spaces-morphisms-lemma-base-change-proper},
\ref{spaces-morphisms-lemma-closed-immersion-proper}, and
\ref{spaces-morphisms-lemma-universally-closed-permanence}).
We have $f_*\mathcal{F} = f_*i_*\mathcal{G} = i'_*(f|_Z)_*\mathcal{G}$
by going around the diagram in two ways. Suppose we know the result holds
for closed immersions and for $f|_Z$. Then we see that
$$
f_*[\mathcal{F}]_k = f_*i_*[\mathcal{G}]_k
= (i')_*(f|_Z)_*[\mathcal{G}]_k =
(i')_*[(f|_Z)_*\mathcal{G}]_k =
[(i')_*(f|_Z)_*\mathcal{G}]_k = [f_*\mathcal{F}]_k
$$
as desired. The case of a closed immersion follows from
Lemma \ref{lemma-length-closed-immersion} and the definitions.
Thus we have reduced to the case where
$\dim_\delta(X) \leq k$ and $f : X \to Y$ is proper and surjective.

\medskip\noindent
Assume $\dim_\delta(X) \leq k$ and $f : X \to Y$ is proper and surjective.
For every irreducible component $Z \subset Y$
with generic point $\eta$ there exists a point $\xi \in X$ such
that $f(\xi) = \eta$. Hence $\delta(\eta) \leq \delta(\xi) \leq k$.
Thus we see that in the expressions
$$
f_*[\mathcal{F}]_k = \sum n_Z[Z],
\quad
\text{and}
\quad
[f_*\mathcal{F}]_k = \sum m_Z[Z].
$$
whenever $n_Z \not = 0$, or $m_Z \not = 0$ the integral closed
subspace $Z$ is actually an irreducible component of $Y$ of
$\delta$-dimension $k$ (see Lemma \ref{lemma-point-of-max-dimension}).
Pick such an integral closed subspace $Z \subset Y$ and denote $\eta$
its generic point.
Note that for any $\xi \in X$ with $f(\xi) = \eta$ we have $\delta(\xi) \geq k$
and hence $\xi$ is a generic point of an irreducible component
of $X$ of $\delta$-dimension $k$ as well
(see Lemma \ref{lemma-point-of-max-dimension}).
By Spaces over Fields, Lemma \ref{spaces-over-fields-lemma-finite-in-codim-1}
there exists an open subspace $\eta \in V \subset Y$
such that $f^{-1}(V) \to V$ is finite.
Since $\eta$ is a generic point of an irreducible component of
$|Y|$ we may assume $V$ is an affine scheme, see
Properties of Spaces, Proposition
\ref{spaces-properties-proposition-locally-quasi-separated-open-dense-scheme}.
Replacing $Y$ by $V$ and $X$ by $f^{-1}(V)$ we reduce to the
case where $Y$ is affine, and $f$ is finite.
In particular $X$ and $Y$ are schemes and we reduce to
the corresponding result for schemes, see
Chow Homology, Lemma \ref{chow-lemma-cycle-push-sheaf}
(applied with $S = Y$).
\end{proof}

\end{document}
